% -------------------------------------------------------------------------
% WBS ID: RES-NUM-FLX
% Deliverable: Spatial Reconstruction and Numerical Flux Selection
% Status: Draft complete
% Evidence status: Regional research specification complete
% Review status: Not reviewed
% Last updated: 2026-07-07
% Dependencies: RES-NUM-FRM, RES-NUM-MSH, RES-MOD-EQS
% Downstream interfaces: RES-NUM-WET, RES-NUM-SRC, RES-VER-BMK, Software flux implementation
% -------------------------------------------------------------------------

\section{RES-NUM-FLX --- Spatial Reconstruction and Numerical Flux Selection}
\label{sec:res-num-flx}

\subsection{Purpose and Scope}
\label{sec:res-num-flx-purpose}

% TODO: Define what this Deliverable covers and explicitly excludes.

\subsection{Research Questions}
\label{sec:res-num-flx-research-questions}

% TODO: State the questions that must be answered before the Deliverable
% can be accepted.

\subsection{Terminology and Definitions}
\label{sec:res-num-flx-definitions}

% TODO: Define the technical terminology, symbols, quantities and
% classifications used in this Deliverable.

\subsection{Literature Evidence}
\label{sec:res-num-flx-literature-evidence}

% TODO: Record evidence from peer-reviewed papers, authoritative datasets,
% standards, technical reports and primary documentation.
%
% Do not create a detached literature-summary list. Explain how each source
% supports, contradicts or limits a project decision.

\subsection{Governing Relations and Quantitative Definitions}
\label{sec:res-num-flx-governing-relations}

% TODO: Record relevant equations, dimensionless groups, empirical models,
% extraction rules or quantitative definitions.
%
% Use align environments for displayed equations.
% Every displayed equation must have a unique \label{eq:...}.
% Do not place punctuation after displayed equations.

\subsection{Discrete Formulation}
\label{sec:res-num-flx-discrete-formulation}

The implemented regional production baseline entering C1A-R10 contains
two hydrodynamic face-state reconstruction options for the same
cell-centred NLSWE finite-volume model. The historical and default path
is first order: each face uses the owner and neighbour cell averages for
$h$, $q_x$ and $q_y$, composed with the existing hydrostatic
well-balanced reconstruction and Rusanov local Lax--Friedrichs flux. The
C1A-R8 global manufactured-solution benchmark verified this path as a
globally first-order approximation of the accepted mathematical model.

The C1A-R9 candidate path is opt-in limited-linear reconstruction. It
reconstructs the primitive hydrodynamic variables
$\eta$, $u$ and $v$, not the conserved momenta directly. For scalar
$\phi\in\{\eta,u,v\}$ the face value from cell $i$ is

\begin{align}
  \phi_{i,f}
  &=
  \phi_i
  +
  \alpha_i^{\phi}
  \nabla\phi_i
  \cdot
  \left(
    \mathbf{x}_f-\mathbf{x}_i
  \right)
  \label{eq:res-num-flx-r9-face-reconstruction}
\end{align}

where $\alpha_i^{\phi}$ is the scalar Barth--Jespersen limiter. The
depth used by the Riemann solve is then recovered from the reconstructed
free surface and the bed elevation. In the verified implementation, the
bed value passed into the hydrostatic reconstruction remains the
cell-local first-order bed value. The higher-order hydrodynamic
reconstruction is therefore composed with, rather than substituted for,
the well-balanced hydrostatic treatment.

Weighted least-squares gradients are formed on the unstructured
cell-centred mesh from face-neighbour samples. For an internal face the
sample is the neighbouring cell centroid and scalar value. For a
boundary face the sample is the boundary face centroid and the already
constructed boundary exterior scalar value, when finite. The weights are
inverse-square-distance weights,
$w=1/\max(\|\mathbf{x}_s-\mathbf{x}_i\|^2,\epsilon_w)$. Cells with
fewer than two valid samples, non-finite data, or an ill-conditioned
$2\times2$ normal system fall back to a zero gradient. This is a local
loss of formal order, not a change in the governing model.

The limiter bounds each reconstructed face value by the local cell and
face-neighbour extrema. It preserves constants, suppresses nonphysical
new extrema at steep gradients, and supports positivity treatment near
wet/dry transitions. The limiter may locally reduce the formal order
near extrema or strongly varying states.

\subsection{Conservation Properties}
\label{sec:res-num-flx-conservation-properties}

% TODO: Complete this RES-NUM Domain-specific subsection.

\subsection{Consistency and Formal Accuracy}
\label{sec:res-num-flx-consistency-formal-accuracy}

% TODO: Complete this RES-NUM Domain-specific subsection.

\subsection{Stability Requirements}
\label{sec:res-num-flx-stability-requirements}

The verified production Regional2D path uses the Rusanov local
Lax--Friedrichs flux with face states rotated into normal/tangential
coordinates. C1A-R9 intentionally did not change the approximate Riemann
solver, SSPRK3 integration, CFL policy, bed-slope physics, Manning
friction, earthquake-source treatment, boundary roles, sponge residuals,
wet/dry handling, coupling variables, or OpenMP residual assembly.

\subsection{Mesh and Time-Step Requirements}
\label{sec:res-num-flx-mesh-time-step-requirements}

% TODO: Complete this RES-NUM Domain-specific subsection.

\subsection{Numerical Failure Modes}
\label{sec:res-num-flx-numerical-failure-modes}

% TODO: Complete this RES-NUM Domain-specific subsection.

\subsection{Verification Requirements}
\label{sec:res-num-flx-verification-requirements}

% TODO: Complete this RES-NUM Domain-specific subsection.

\subsection{Candidate Approaches}
\label{sec:res-num-flx-candidate-approaches}

% TODO: Define the credible candidate approaches considered.

\subsection{Critical Comparison}
\label{sec:res-num-flx-critical-comparison}

% TODO: Compare candidates using physically and project-relevant criteria.
% Do not select an approach solely on popularity or implementation ease.

\subsection{Assumptions and Applicability}
\label{sec:res-num-flx-assumptions}

% TODO: State assumptions and the conditions under which they are valid.

\subsection{Limitations and Failure Conditions}
\label{sec:res-num-flx-limitations}

% TODO: State where the evidence, model or selected approach ceases to be
% reliable or applicable.

\subsection{Project-Specific Conclusions}
\label{sec:res-num-flx-conclusions}

% TODO: Convert the evidence into conclusions specific to the Studentship
% project. Do not merely restate literature findings.

\subsection{Selected Approach and Rationale}
\label{sec:res-num-flx-selected-approach}

% TODO: Record the selected approach only after sufficient evidence exists.
% State the alternatives rejected and the reasons for rejection.

Selected regional flux/reconstruction hierarchy for the verified
implementation:

\begin{itemize}
  \item first-order piecewise-constant hydrodynamic reconstruction as
    the default and reference path;
  \item opt-in limited-linear reconstruction in smooth wet regions;
  \item weighted least-squares gradients on the face-neighbour stencil;
  \item scalar Barth--Jespersen limiting;
  \item first-order bed values inside the existing hydrostatic
    reconstruction;
  \item Rusanov local Lax--Friedrichs numerical flux;
  \item conservative face-based residual assembly and SSPRK3 transport.
\end{itemize}

\subsection{Unresolved Decisions}
\label{sec:res-num-flx-unresolved-decisions}

% TODO: Record unresolved questions, required evidence and decision owners.

The first-order and limited-linear reconstruction paths are no longer
open for controlled verification: C1A-R8 verified the former as globally
first order and C1A-R9 verified the latter as approximately second order
on smooth MMS benchmarks. Their event-level behaviour in the
Tōhoku--Kamaishi forcing corridor remains the R10 numerical question.
Alternative flux families such as HLLC remain deferred research options,
not the verified R10 production candidate.

\subsection{Requirements and Handoffs}
\label{sec:res-num-flx-requirements}

% TODO: State requirements passed to other Research Domains, Software,
% verification activities or later execution work.

\subsection{Deliverable Completion Record}
\label{sec:res-num-flx-completion-record}

% TODO: Record whether the Deliverable is:
%
% - Not started
% - In progress
% - Draft complete
% - Reviewed
% - Accepted
%
% Acceptance requires:
%
% - the research questions to be answered;
% - claims to be supported by appropriate evidence;
% - assumptions and limitations to be explicit;
% - the project-specific conclusion to be stated;
% - downstream requirements to be traceable;
% - unresolved decisions to be closed or formally deferred.
